\subsection{实验结果分析}
\begin{itemize}
    \item 基线版本与优化一分析:
    \begin{itemize}
        \item 现象:基线版本的周期数为\texttt{50012}，优化一的周期数下降到\texttt{40006}，在保证结果仍然为\texttt{50005000}的前提下，性能得到了明显提升。
        \item 原因:基线版本在循环中需要频繁调用\texttt{vec\_length(v)}和\texttt{get\_vec\_element()}，同时还存在较多的边界检查和不必要的中间访存开销。优化一通过直接访问\texttt{v->len}和\texttt{v->data}，并使用局部变量完成累加后统一写回，减少了过程调用和存储访问次数，因此执行效率明显提高。
    \end{itemize}

    \item 优化二分析:
    \begin{itemize}
        \item 现象:在优化一基础上引入4路循环展开后，周期数进一步下降到\texttt{25032}，相比优化一有了更加显著的提升，说明该优化在当前实验环境下效果十分明显。
        \item 原因:4路循环展开使每轮循环能够同时处理4个元素，从而减少了循环控制指令（如循环判断、跳转、计数更新）所占比例。同时设置多个独立累加器后，可以减轻单一累加变量带来的顺序相关问题，使程序更容易发挥指令级并行能力，因此整体性能得到了大幅改善。
    \end{itemize}

    \item 优化三分析:
    \begin{itemize}
        \item 现象:优化三在优化二基础上改用指针遍历并保留4路展开后，周期数变为\texttt{27536}，反而略高于优化二的\texttt{25032}，说明该优化没有继续带来性能提升。
        \item 原因:虽然指针遍历在理论上可能减少部分地址计算开销，但在当前RISC-V裸机仿真环境与编译条件下，这种写法并未优于数组下标访问。相反，优化三引入的额外控制逻辑可能抵消了部分收益，因此其实际运行效果不如优化二。这也说明并不是所有“更底层”的写法都一定更快，优化效果仍需通过实验数据验证。
    \end{itemize}

    \item 优化四分析:
    \begin{itemize}
        \item 现象:优化四在优化二思路基础上进一步扩展为8路累加器和8路循环展开，周期数继续下降到\texttt{22556}，成为所有测试版本中的最佳结果。
        \item 原因:优化四相比优化二进一步扩大了循环展开规模，并增加了独立累加器数量，使循环控制开销进一步减少，同时也进一步降低了顺序相关对流水线执行效率的限制。因此，该版本在当前实验环境下能够获得最好的性能表现。
    \end{itemize}

    \item 编译器优化尝试分析:
    \begin{itemize}
        \item 现象:在优化四基础上增加函数级\texttt{O3}编译器优化属性后，程序输出结果仍然正确，但周期数与优化四相比没有发生变化，说明未观察到额外明显收益。
        \item 原因:这说明优化四本身已经较充分地利用了当前编译器可发挥的优化空间，其代码结构已经较接近当前实验环境下的高效实现形式。因此，继续显式增加\texttt{O3}并没有带来进一步的可观提升。本实验中的主要性能收益仍然来自手工代码优化，而不是额外加入的编译器优化属性。
    \end{itemize}

    \item 总体结果分析:
    \begin{itemize}
        \item 现象:从基线版本到优化四，周期数由\texttt{50012}逐步下降到\texttt{22556}，说明总体优化过程是有效的；其中优化一、优化二和优化四均取得了明显提升，而优化三和编译器优化尝试未带来进一步改善。
        \item 原因:本实验中性能提升的主要来源包括减少过程调用、降低不必要的存储引用、采用循环展开以及增加独立累加器数量等手工代码优化方法。其中，优化四在减少循环控制开销和提高指令级并行性方面效果最好，因此成为最终性能最优的实现版本。
    \end{itemize}

    \item 最终结论:
    \begin{itemize}
        \item 若以基线版本作为参照，则优化四相对基线的性能提升约为
        \[
        \frac{50012 - 22556}{50012} \approx 54.9\%
        \]
        \item 这表明，在保证程序计算结果正确的前提下，优化四将核心求和过程的执行周期降低了约\textbf{54.9\%}，是本实验中最有效的优化方案。
    \end{itemize}

\end{itemize}